\chapter[Pacman and neutral bounds]{Source digests IV: Pacman renormalization and neutral bounds}
\label{ch:digests-pacman-neutral}

This chapter digests two recent papers that the survey~\cite{DudkoMLC} uses as
building blocks: the Pacman renormalization near Siegel parameters of
Dudko, Lyubich and Selinger~\cite{DudkoLyubichSelinger}, and the uniform a priori
bounds for neutral renormalization of Dudko and Lyubich~\cite{DudkoLyubichNeutral}.
Statements are those of the sources, up to notation; theorem numbers are those
printed in the sources.

\section{Pacman renormalization near Siegel parameters}
\label{sec:digest-pacman}

The paper~\cite{DudkoLyubichSelinger} places the Branner--Douady surgery relating
limbs of $\M$ in the framework of Pacman renormalization.  Its authors are
Dudko, Lyubich and Selinger.  (The survey cites this paper as its reference~[20]
for the Pacman/Siegel hyperbolicity result.)

\subsection*{Definitions}

\begin{definition}[Notation for rotations]
\label{def:pacman-rotation}
Put $e(z)=e^{2\pi iz}$ and $p_\theta(z)=e(\theta)z+z^2$.  The set $\Theta_{\rm bnd}$ consists of
rotation numbers with bounded continued-fraction entries, and $\Theta_{\rm per}$ of
numbers periodic under the map
\[
 \mathcal R(\theta)=\begin{cases}\theta/(1-\theta), & 0\le\theta\le\tfrac12,\\
                      (2\theta-1)/\theta, & \tfrac12\le\theta\le1.\end{cases}
\]
A Siegel map has a neutral fixed point with multiplier $e(\theta)$, a quasidisk
Siegel disk compactly contained in its domain, and exactly one critical point on the
boundary of the Siegel disk.
\end{definition}

\begin{definition}[Full pacman; truncated pacman]
\label{def:pacman-full}
A full pacman is a map $f\colon U\to V$, where $V$ is a closed topological disk with a
critical arc $\gamma_1$ joining $\partial V$ to a fixed point $\alpha\in\operatorname{int}V$,
such that $U\Subset V$ and $\gamma_1$ has exactly three lifts: $\gamma_0\subset U$ starting
at $\alpha$, and $\gamma_\pm\subset\partial U$ starting at the pre-fixed point $\alpha_0$.
Moreover $f\colon U\setminus\gamma_0\to V\setminus\gamma_1$ is a two-to-one branched covering,
and $f$ extends locally conformally across $\partial U\setminus\{\alpha_0\}$.  A truncated
pacman is obtained by removing a small disk around $\alpha_0$.
\end{definition}

\begin{definition}[Dynamical objects of a pacman]
\label{def:pacman-dynamics}
The non-escaping set is $K_f=\bigcap_{n\ge0}f^{-n}(U)$, and $V\setminus K_f$ is the escaping
set.  The pacman Julia set is $J_f=\bigcup_{n\ge0}f^{-n}(\partial Z_f)$, where $Z_f$ is the
Siegel disk when there is one.  External rays are leaves of the pulled-back lamination.
\end{definition}

\begin{definition}[Prepacman; pacman renormalization]
\label{def:pacman-renorm}
A prepacman is a sector $S$, split by an interior ray into two pieces $T_\pm$, with
holomorphic maps $f_\pm\colon U_\pm\to S$, $U_\pm\subset T_\pm$, such that a gluing map
$\psi\colon S\to V$ projects $(f_-,f_+)$ to a pacman.  A map $f$ is renormalizable if a
prepacman $G=(f^a\colon U_-\to S,\ f^b\colon U_+\to S)$ realizes the first return to $S$,
with the preceding orbits covering a neighbourhood of $\alpha$ compactly contained in $U$.
The glued pacman is the renormalization $Rf$.  The return times $a,b$ are recorded, and the
renormalization is prime when $a+b=3$.
\end{definition}

\begin{definition}[Siegel pacman; standard Siegel pacman]
\label{def:pacman-siegel}
A pacman is Siegel if it has a Siegel disk $Z_f$, $\gamma_1$ is an external ray followed by
an internal ray that meets $\partial Z_f$ at a non-precritical point, and the truncation disk
is an equipotential subdisk of $Z_f$.  A standard Siegel pacman is one in which $\gamma_0$
passes through the critical value, and the standard renormalization is based at the
critical value.
\end{definition}

\begin{definition}[Satellite valuable flower]
\label{def:pacman-flower}
An open forward-invariant set $H$ with $H\cup\{\alpha\}$ connected, divided into $q$
petals $H_i$ with $g(H_i)\subset H_{i+p}$, such that an attracting $q$-cycle attracts all
points of $H$ and $H_{-p}$ contains the critical point.  Its combinatorial rotation number is
$p/q$, and its multiplier is that of the attracting cycle.
\end{definition}

\subsection*{Results}

\begin{theorem}[Dudko--Lyubich--Selinger, Theorems~1.1 and~7.7]
\label{thm:dls-1-1}
For every $\theta\in\Theta_{\rm per}$ the renormalization $R$ has a unique periodic Siegel
pacman $f_*$.  It is hyperbolic with one-dimensional unstable manifold.  Locally, its stable
manifold is the set of pacmen with the same multiplier, and every point of it is Siegel.
Status: proved in the paper.
\end{theorem}

\begin{theorem}[Dudko--Lyubich--Selinger, Corollaries~1.2 and~7.9]
\label{thm:dls-1-2}
Nearby maps with a fixed rotation number $\theta\in\Theta_{\rm per}$ have Siegel disks that
depend continuously on the map, and nearby locally Siegel maps remain Siegel.  Status:
proved in the paper, using earlier results.
\end{theorem}

\begin{theorem}[Dudko--Lyubich--Selinger, Theorem~1.3]
\label{thm:dls-1-3}
If $p_n/q_n$ are the continued-fraction approximants of $\theta\in\Theta_{\rm per}$, then
$|c(\theta)-a_{p_n/q_n}|\sim q_n^{-2}$.  Status: proved in the paper via Theorem~8.2.
\end{theorem}

\begin{theorem}[Dudko--Lyubich--Selinger, Theorem~3.2; bounded type]
\label{thm:dls-3-2}
For $\theta\in\Theta_{\rm bnd}$ the closed Siegel disk of $p_\theta$ is a quasidisk containing
the critical point.  Conversely, a one-critical-point map whose compactly contained Siegel
disk is a quasidisk containing the critical point has a bounded-type rotation number.  Status:
the paper uses results of Douady--Ghys, Herman, Świątek and Avila--Lyubich.
\end{theorem}

\begin{theorem}[Petersen, Theorem~3.3 of the paper; cited]
\label{thm:dls-petersen}
For the bounded-type rotation numbers of the paper, the filled Julia set $K_{p_\theta}$ is locally connected,
and the components that remain after removing bubbles of bounded generation have arbitrarily small
diameter.  Status: cited from Petersen.
\end{theorem}

\begin{theorem}[Dudko--Lyubich--Selinger, Theorem~3.12]
\label{thm:dls-3-12}
For every standard Siegel pacman, $K_f$ is locally connected; the components of bounded
generation shrink, and every periodic point of $J_f$ is the landing point of a bubble chain.
Status: proved in the paper.
\end{theorem}

\begin{theorem}[Dudko--Lyubich--Selinger, Theorem~3.16]
\label{thm:dls-3-16}
Near $f_*$ the renormalization $R$ is an analytic compact operator, and $Rf_*=f_*$.  Status:
proved in the paper.
\end{theorem}

\begin{theorem}[Dudko--Lyubich--Selinger, Theorem~8.2; satellite flowers]
\label{thm:dls-8-2}
For a sequence $p_n/q_n\to\theta_*$ satisfying the displayed renormalization relation, a
transverse one-dimensional slice contains continuous paths of satellite valuable flowers with
rotation numbers $p_n/q_n$ and a prescribed multiplier path.  The exponent in the displayed
distance estimate is not legible in the plain-text version examined, so it is not recorded here.
Status: proved in the paper.
\end{theorem}

\subsection*{Relation to MLC}

The paper proves local connectivity of $K_f$ for standard Siegel pacmen
(Theorem~\ref{thm:dls-3-12}) and recalls the local connectivity of Julia sets of
bounded-type Siegel polynomials (Petersen).  It does not prove local connectivity of
$\M$ itself.  It states that a later work applies these methods to some bounded-type
satellite parameters.  Its appendix states that a molecule-conjecture framework would
imply MLC for all infinitely renormalizable satellite-type parameters; this is conditional
and is not claimed here.

\section{Uniform a priori bounds for neutral renormalization}
\label{sec:digest-neutral-apriori}

The paper~\cite{DudkoLyubichNeutral} proves uniform bounds for neutral quadratic polynomials
$f_\theta(z)=e(\theta)z+z^2$, with $\theta$ in the bounded-type set.

\subsection*{Definitions}

\begin{definition}[Critical point and combinatorial intervals]
\label{def:neutral-critical}
For $f_\theta$ the critical point is $c_0(f_\theta)=-e(\theta)/2$.  If $p_n/q_n$ are the best
approximants of $\theta$ and $Z$ is the Siegel disk, the level-$n$ interval is
$I^n_f(x)=[x,f^{q_n}(x)]\subset\partial Z$.  Its combinatorial length is $|I|=|h(I)|_{\R/\Z}$,
where $h\colon(Z,\alpha)\to(\mathbb D,0)$ is the linearizing map.
\end{definition}

\begin{definition}[Pseudo-Siegel bounds]
\label{def:neutral-pseudo-siegel}
$f_\theta$ has pseudo-Siegel bounds if there are disks
$\widehat Z^{-1}_\theta\supset\widehat Z^0_\theta\supset\widehat Z^1_\theta\supset\cdots$ with
$\bigcap_{m\ge-1}\widehat Z^m_\theta=H_\theta$, such that $\widehat Z^m_\theta$ is almost invariant
under $f_\theta^{q_{m+1}}$, its boundary carries a nest of tilings of essentially bounded
geometry, and $f_\theta$ is injective on the relevant disk.  The bounds are uniform when the
constants of the tilings are independent of $\theta$.
\end{definition}

\begin{definition}[Mother Hedgehog]
\label{def:neutral-hedgehog}
A hull is a compact connected full set.  A Mother Hedgehog of $f_\theta$ is an invariant hull
$H_f$ containing $0$ and $c_0(f)$.
\end{definition}

\begin{definition}[Parameter compactification]
\label{def:neutral-theta}
Put $\Theta=\{\theta=[0;a_1,a_2,\dots]:\ a_i\in\N_{\ge1}\cup\{\infty\}\}$ and
$\underline\Theta=(\R\setminus\Q)/\Z$.  The boundary of $\Theta$ corresponds to Lavaurs--Epstein
towers.
\end{definition}

\begin{definition}[$\delta$-pseudo-Siegel disk]
\label{def:neutral-delta}
Inductively, $\widehat Z^n=Z$ for large $n$, and either $\widehat Z^m=\widehat Z^{m+1}$, or
$\widehat Z^m=\widehat Z^{m+1}\cup Z^m$ where $Z^m$ is a $\delta/2$-near-rotation domain called the
core.  The second operation is the regularization at level $m$; it fills sufficiently deep
parabolic fjords.
\end{definition}

\begin{remark}
The paper refers to $\psi^\bullet$-quasi-quadratic-like maps and to neutral renormalization
towers, but defines neither in the present text; it refers both to an external paper cited as
{[DLL]}.
\end{remark}

\subsection*{Results}

\begin{theorem}[Dudko--Lyubich, Theorem~1.1; uniform bounds]
\label{thm:neutral-1-1}
There is an absolute constant $K$ such that, for every bounded-type Siegel quadratic polynomial
$f=f_\theta$ and every combinatorial interval $I=I^n_f(x)$, one has $W^+_3(I)\le K$.  Status:
proved in the paper.
\end{theorem}

\begin{theorem}[Dudko--Lyubich, Theorem~1.2; Mother Hedgehog]
\label{thm:neutral-1-2}
For every neutral quadratic polynomial $f=f_\theta$ with $\theta\notin\Q$ there is a Mother
Hedgehog $H_f\ni c_0(f)$ such that $f\colon H_f\to H_f$ is a homeomorphism.  Status: proved in the
paper (also Theorem~11.2).
\end{theorem}

\begin{theorem}[Dudko--Lyubich, Theorem~1.3; quasidisk approximation]
\label{thm:neutral-1-3}
There is an absolute $K$ such that every bounded-type Siegel quadratic polynomial has a disk
$\widehat Z_f\supset Z_f$ which is a $K$-quasidisk and on which $f$ is injective.  Status: proved in
the paper (detailed as Theorem~11.1).
\end{theorem}

\begin{theorem}[Sector bounds; cited in the paper as {[DL2]}]
\label{thm:neutral-1-5}
There is a compact sector renormalization operator $R_{\rm sec}$ for all neutral quadratic polynomials,
with the properties (I)--(VI) listed in the paper: uniform continuity of the pseudo-Siegel disks,
existence throughout $\Theta$, continuous Mother Hedgehogs, precompact transcendental closest-return
limits, compatibility with renormalization, and star-like Mother Hedgehogs.  Status: cited in the
paper as {[DL2]}; the reference list of the paper gives its details.
\end{theorem}

\begin{theorem}[QL Siegel bounds; cited in the paper as {[DLL]}]
\label{thm:neutral-1-6}
If $g\colon X\to Y$ is quadratic-like with bounded-type Siegel disk $Z_g$, and $K_g=W_\bullet(g)=W(Y\setminus Z_g)$,
then for $I=[x,g^{q_{n+1}}x]$ one has $W^+_3(I)=O(K_g/q_{n+1})$.  In particular bounded $K_g$ gives
$W^+_3(I)=O(1)$.  Status: cited in the paper as {[DLL]}.
\end{theorem}

\begin{theorem}[Dudko--Lyubich, Theorem~10.1; golden mean]
\label{thm:neutral-10-1}
For an eventually golden-mean $\theta$ there is an absolute $K$ with $W^+_3(I)\le K$ for every combinatorial
interval $I\subset\partial Z$.  Status: proved in the paper.
\end{theorem}

\begin{theorem}[Dudko--Lyubich, Theorem~10.2; scale dichotomy]
\label{thm:neutral-10-2}
For the same $\theta$ there are absolute $N,K$ and a sequence of geodesic pseudo-Siegel disks such that, if
$\ell_{m+1}\le|I|\le\ell_m$, then either $|I|\ge\ell_m/2$, or $|I|\le N\ell_{m+1}$ and $W^+_3(I)\le K$.
Status: proved in the paper.
\end{theorem}

\subsection*{Relation to MLC}

The paper states that Yoccoz proved MLC at parameters that are not infinitely quadratic-like
renormalizable, and it formulates the principle that a priori bounds for every combinatorial class
should give quasiconformal conjugacy, hence singleton parameter classes, hence MLC.  It notes that
quadratic-like bounds in general fail in the satellite case when the Julia set is not locally connected
at the critical point.  Its authors state that the uniform neutral bounds control satellite
renormalization relevant to MLC, and they conjecture that pseudo-Siegel bounds for satellite
renormalization, interpolated with the primitive bounds of Kahn--Lyubich, would complete the proof
of MLC.  that conjecture is not proved here.  The paper also refers to MLC at bounded-type parameters
proved in an earlier work by the same authors, cited there under its own label.